\section{Persistence, temporary shocks, and lock-in}
\label{app:persistence}

This appendix collects the results that start from an equilibrium and explicit
basin conditions: the general existence route, the conditional high-basin
theorem and its long run, dual-basin path dependence, the effects of temporary
sanction and suspensions of sanctions, the comparative statics of lock-in, and the
contracting closed loop of the class under the lifting condition. \Cref{sec:discussion}
summarizes what they establish. The statements are those of the previous main
text, moved without change except the one correction noted at
\Cref{cor:pathwise-dominance}, the restatements of 2026-10-07, which keep
each conclusion and drop hypotheses the proofs do not read, and the
2026-10-08 restatement of \Cref{cor:pathwise-dominance} to the hypotheses and
existence conclusion of its certificate theorem; A1--A3 and the
basin, which only these results use, come first.

\subsection{Assumptions and the basin}
\label{sec:assumptions}

The high-basin and pathwise-ordering results of this appendix require only
\textup{(F2)--(F4)} of the transition; the shock-entry result
(\Cref{thm:shock}) additionally uses continuity \textup{(F1)} and full-sanction
ascent below the basin, which \textup{(F5)} implies. For the optional
quantitative value-regularity route, we additionally use an inherited-state
modulus $\gamma_F$:
\begingroup
\renewcommand{\theHequation}{F6Q}
\begin{equation}
\bigl|\F(x_2,s)-\F(x_1,s)\bigr|
\leq \gamma_F |x_2-x_1|
\quad\text{for all }x_1,x_2,s\in[0,1],
\qquad
\beta_S\gamma_F<1.
\tag{F6-Q}\label{eq:F-state-lip}
\end{equation}
\endgroup
Condition \textup{(F6-Q)} is not part of the qualitative baseline
\textup{(F1)--(F5)}, and it is stated on the feasible square. Under
\eqref{eq:affine}, $\gamma_F=1-\delta$: it is the coefficient on the inherited
state, not the sanction pass-through $\delta$ (\Cref{app:moved-ah}).

We state three structural inputs and derive the high political basin.
A1 and A2 restrict objectives uniformly over value sets; A3 is a
curvature condition on the sender's stage payoff. The general existence
theorem takes its invariant-set, selector and operator-bound conditions as
hypotheses; the constructions of \Cref{sec:existence:placable} and
\Cref{sec:existence:nonvacuity} build their invariant sets and regularity
constants from $(U_T,U_S,\F,\beta_T,\beta_S)$ in closed form in the stated rates
(\Cref{app:closedforms}). The derived-basin
route in \Cref{sec:flip} derives the basin and the target condition from
complementarity. The three assumptions also provide a fallback for
specifications without global complementarity.

\subsubsection*{Basin pivots}

Let $\bar x_T,\bar x_S\in[0,1]$ denote the target and sender tipping
thresholds. Under stage complementarity, each is the infimum of a flip
set; \Cref{prop:tipping} establishes existence, interiority, and the
upper-interval structure. Define the joint pivot
$\bar x=\max\{\bar x_T,\bar x_S\}$ and the high basin
$\basin=[\bar x,1]$.

\subsubsection*{Lipschitz constants}

Write $K_S$ for a Lipschitz constant, in the inherited state, of the
sender's values in its value set. Under the optional regularity \eqref{eq:F-state-lip} the
envelope recursion gives $K_S=L_{S,x}^{\alpha}/(1-\beta_S\gamma_F)$, where
$L_{S,x}^{\alpha}$ bounds the state variation of the sender's stage payoff along
a continuous target policy \citep{StokeyLucas1989,BertsekasTsitsiklis1996};
the target constant and the selector moduli are in \Cref{app:closedforms}.
The modulus $\gamma_F$ does not bound the transition's response to the
sanction, so the level condition used below is the transition-adjusted slope condition
\begin{equation}
\Us(x,a,s_2)-\Us(x,a,s_1)
>
\beta_S K_S\bigl|\F(x,s_2)-\F(x,s_1)\bigr|,
\qquad s_1<s_2,
\label{eq:F-adjusted-sender-chord}
\end{equation}
on the range used below. It needs no separate modulus of $\F$ in the sanction
and, under the partial-adjustment law, follows from $\beta_SK_S\delta<m_S$ for a stage
lower slope $m_S$ (\Cref{app:moved-ah}).

\subsubsection*{Three structural assumptions}

We refer to \Cref{ass:sender,ass:target,ass:concavity} as A1--A3. A1
controls the direction of the sender's response, A2 records the
target's forward-looking preference on the basin, and A3 gives
diminishing returns in the sender's stage problem. The high-basin
theorem uses the direction and target conditions together with a separate
sanction-level condition, the target condition only for its conclusion (vi);
the shock theorem uses no target condition, and neither uses A3 directly.
On the derived-basin route A2 can be derived rather than imposed, though the
route's conclusions no longer read it (\Cref{prop:flip-route}). \Cref{tab:assumption-routes}, in \Cref{app:moved-ad},
records which
result consumes which of these conditions: none is a maintained assumption
on the whole model, and the constructed equilibria of \Cref{sec:four-regimes}
use none of A1--A3 as stated.


Write $\mathcal C_T$ and $\mathcal C_S$ for the target and sender value
sets: closed, nonempty sets of bounded continuous functions on
$[0,1]$ cut out by a sup-norm radius together with Lipschitz or slope bounds,
with radii and moduli derived from the model primitives in
\Cref{sec:existence,app:closedforms}. The general existence theorem takes these
sets and their invariance as inputs; the constructions discharge
them, the target's invariant set on the class under the lifting condition by a sup-norm radius and a lower-slope
condition with no Lipschitz clause.

\begin{assumption}[Sender interval dominance after target-policy substitution]
\label{ass:sender}
For every value pair
$(\Vt,\Vs)\in\mathcal C_T\times\mathcal C_S$, let
$\alpha_{\Vt,\Vs}$ be its induced target policy and define the
policy-substituted sender family by
\begin{equation}
\widetilde J_{\Vt,\Vs}(s,x)
:=
J_S\bigl(x,\alpha_{\Vt,\Vs}(x),s;\Vs\bigr).
\label{eq:sender-policy-substituted}
\end{equation}
On the sanction--state domain $[0,1]\times\basin$, this family
satisfies the interval dominance order in $(s,x)$ of
\citet{QuahStrulovici2009}. Milgrom--Shannon single crossing in
$(s,x)$ is a sufficient, stronger condition.
\end{assumption}

A1 is ordinal: after the endogenous target policy is substituted, it makes the
sender's dynamic best response non-decreasing in the state; single crossing is a
stronger sufficient condition \citep{MilgromShannon1994}. Direction does not
determine level: keeping the response above $\max\{x,\bar x\}$ is the role of
\eqref{eq:F-adjusted-sender-chord}, and the proof of \Cref{thm:persistence} keeps the two
arguments separate (\Cref{app:proofs,app:dischargers,app:moved-ah}).

\begin{assumption}[Target combined slope condition, value-set form]
\label{ass:target}
For every $(V_T,V_S)\in\mathcal C_T\times\mathcal C_S$, let
$\alpha_{V_T,V_S}$ be the induced target policy. For every
$x\in\basin$,
\begin{equation}
\boxed{
 s\longmapsto
 U_T\bigl(x,\alpha_{V_T,V_S}(x),s\bigr)
 +\beta_TV_T\bigl(\F(x,s)\bigr)
 \text{ is non-decreasing on }[\max\{x,\bar x\},1].
}
\label{eq:A2}
\end{equation}
\end{assumption}

A2 is imposed uniformly over the value sets rather than at a presumed
equilibrium. It combines the target's current payoff from sanctions with
the political-coalition effect in its continuation value
\citep{Kaempfer1988}. Each component may be non-monotone; their sum must
be non-decreasing on the displayed range. Strict concavity of the
complete target objective in compliance, together with Berge's Maximum
Theorem, supplies a continuous policy selector
\citep{AliprantisBorder2006}.

\begin{assumption}[Uniform stage strong concavity in $s$]
\label{ass:concavity}
There is $\mu>0$ such that, for every $(x,a)\in[0,1]^2$, every
$s_1,s_2\in[0,1]$ and every $t\in[0,1]$,
\[
t\,U_S(x,a,s_1)+(1-t)\,U_S(x,a,s_2)
\le U_S\bigl(x,a,ts_1+(1-t)s_2\bigr)-\tfrac{\mu}{2}\,t(1-t)(s_1-s_2)^2 .
\]
This is the derivative-free form the formalization consumes. When $U_S$ is twice
differentiable in $s$ it is the familiar
$\sup_{(x,a,s)\in[0,1]^3}\partial^2 U_S/\partial s^2\le-\mu<0$.
\end{assumption}

A3 imposes diminishing returns to sanctions, as in standard public-choice
models \citep{Kaempfer1988,Eaton1992}, and makes the myopic best response
unique. It does not make the maximizer of $J_S$ unique, because the continuation
term can offset stage curvature, and it is not what carries the substitution
between compliance and the sanction (\Cref{app:closedforms,app:moved-ah}).

\paragraph{A route that derives the basin.}
Under a stage-complementarity assumption the pivots $\bar x_T$ and $\bar x_S$
are not inputs: each is the infimum of a payoff flip set that is an upper
interval, and the target condition A2 becomes a result
(\Cref{ass:flip,prop:tipping,prop:flip-route}, in \Cref{sec:flip}). Sign
conditions of this kind fix the signs of slopes and bound no magnitude, so they do not
deliver persistence on their own; the magnitude condition is the sub-pivot
slope condition (P5) of \Cref{thm:persistence}. No result on the class under the lifting condition in
\Cref{sec:four-regimes} uses this route.

\subsection{Equilibrium existence and persistence in the general game}
\label{sec:existence:general}

The dynamic results of this subsection require a stationary feedback
Stackelberg equilibrium, and they start from one.

\subsubsection{Equilibrium existence}

A joint contraction gives one route to existence in a smooth,
gently coupled region. Let $L_{TS}$ and $L_{ST}$ measure how each
Bellman row responds to the other player's continuation value, and let
$\beta_{S,\mathrm{row}}$ be the sender row's own contraction modulus;
$\rho(\mathbf M)$ is the spectral radius of the coupling matrix these
four numbers form (\Cref{app:existencemachinery}).

\Cref{thm:existence} (\Cref{app:existencemachinery}) states the route and its
Schur condition $L_{TS}L_{ST}<(1-\beta_T)(1-\beta_{S,\mathrm{row}})$, with no
sign condition on the four numbers. Its uniqueness conclusion --- every
equilibrium with values in the invariant sets that plays the target selector those sets define
coincides with the constructed one --- is confined to the stated invariant-set and
selector restrictions; it is not global equilibrium uniqueness. A summary of the
statement is in \Cref{app:moved-ad}.

The economically transparent deterministic route is the high-basin
corner. When full non-compliance is optimal for the target, the game reduces
to sequential single-agent Bellman problems.

\begin{theorem}[Basin-corner feedback Stackelberg existence]
\label{thm:corner-existence}
Fix $\Vs$ at its Bellman fixed point against $\alpha\equiv0$, and
write $\sigma^\ast=\sigma^{\mathrm{dyn,top}}(\Vs)$. Suppose, on the
entire feasible state space, that:
\begin{enumerate}[label=(\roman*)]
\item the map $(x,a)\mapsto\sigma^\ast(x,a)$ is continuous at feasible
states $x$ and every target action $a$;
\item the target continuation value constructed against
$\sigma^\ast$ is non-decreasing in $x$;
\item for $a\in(0,1]$, $\sigma^\ast(x,a)\le\sigma^\ast(x,0)$: no positive
compliance raises the sender's response above its response to full
non-compliance; and
\item the substituted target stage payoff is highest at full non-compliance,
\[
U_T\!\bigl(x,a,\sigma^\ast(x,a)\bigr)
\le
U_T\!\bigl(x,0,\sigma^\ast(x,0)\bigr)
\qquad\text{for }a\in(0,1].
\]
\end{enumerate}
Then $\alpha\equiv0$ is dynamically optimal and
$E=(\Vt,\Vs,\alpha\equiv0)\in\mathcal E$, with $\Vt$ constructed as
the single-agent Banach fixed point.
\end{theorem}

Continuity gives a well-defined target Bellman operator, while the
remaining conditions make full non-compliance optimal after the sender's
response is substituted. Each condition is used once and only weakly: an
antitone response and a strict stage decrease, the form stated before
2026-10-07, are a special case. This argument applies to the general
monotone transition law and does not require purification.

Off the corner the deterministic joint iteration need not preserve order or
contract; existence there is obtained for a transition-smoothed game (an atomless
disturbance added to the state law; ``purified'' in earlier versions, a term this
literature reserves for payoff noise) and at example economies, not for
the primitive class (\Cref{app:moved-ad,app:scope-notes}).

The dynamic results use the general transition law
\eqref{eq:transition}. Monotonicity and the diagonal condition deliver
basin invariance; strict adjustment makes the long-run sanction
equal to the limiting state. The partial-adjustment law \eqref{eq:affine} is needed only for the
closed-form expressions stated separately below.

\subsection{Persistence and the long run}

\begin{theorem}[Persistence on the high political basin]
\label{thm:persistence}
Let $\bar x_T,\bar x_S\in[0,1]$ and
$\bar x=\max\{\bar x_T,\bar x_S\}$. Fix the value sets
$\mathcal C_T\times\mathcal C_S$ and let
$E=(\Vt,\Vs,\alpha)\in\mathcal E$. Suppose:
\begin{enumerate}[label=(P\arabic*),leftmargin=2.7em]
\item $\Vt\in\mathcal C_T$, $\Vs\in\mathcal C_S$, and
$\alpha=\alpha_{\Vt,\Vs}$;
\item every sender value in $\mathcal C_S$ is $K_S$-Lipschitz;
\item A1 holds for every value pair, in the state-indexed form
\eqref{eq:sender-policy-substituted};
\item the target combined slope condition A2 holds for every value pair; and
\item for every value pair $(V_T,V_S)$, every basin state $x\in\basin$,
and every sanction $0\le s<x$ below the pivot $\max\{x,\bar x\}=x$, the
general-$\F$ sub-pivot slope condition holds in pivot-star form:
\[
\Us\!\left(x,\alpha_{V_T,V_S}(x),x\right)
-
\Us\!\left(x,\alpha_{V_T,V_S}(x),s\right)
>
\beta_S K_S\bigl|\F(x,x)-\F(x,s)\bigr|.
\]
\end{enumerate}
Then for every $x_0\in\basin$, the equilibrium path satisfies:
\begin{enumerate}[label=(\roman*)]
\item $x^E_t\in\basin$ for every $t\ge0$;
\item $s^E_t\ge\bar x$ for every $t\ge0$;
\item $s^E_t\le s^E_{t+1}$ for every $t\ge0$;
\item $\alpha(x_t^E)$ is feasible and maximizes the target's complete
Bellman objective against the sender's dynamic best response;
\item $s_t^E$ maximizes $J_S(x_t^E,\alpha(x_t^E),\cdot\,;\Vs)$; and
\item the map
\[
s\longmapsto
\Ut(x_t^E,\alpha(x_t^E),s)
+\beta_T\Vt\bigl(\F(x_t^E,s)\bigr)
\]
is non-decreasing on $[\max\{x_t^E,\bar x\},1]$.
\end{enumerate}
\end{theorem}

The proof separates the direction and the level of the sender's response: A1
makes the dynamic best response non-decreasing in the inherited state, and the
transition-adjusted slope condition P5 keeps it above the basin floor, so that
$s_t^E\ge x_t^E$ and the basin is preserved. P5 is read only at basin states
and only against the pivot: it compares each sub-pivot sanction with the pivot,
never two sub-pivot sanctions with each other, and it constrains no state below
the joint pivot; the pairwise slope condition on $[\bar x_S,1]$ implies it. The target
condition P4 is logically
separate and yields only conclusion (vi); conclusions (i)--(iii) hold under
P1--P3 and P5 alone. A1 is used only for (iii): conclusions (i), (ii) and
(iv), and a non-decreasing state path, follow from P1, P2 and P5, the level
condition, without it. The theorem is conditional on the stated equilibrium and
its value-set conditions; \Cref{thm:corner-existence} and
\Cref{prop:primitive-persistence} provide unconditional applications under
narrower primitive conditions. The proof is in \Cref{app:proofs}; a local form
with no value sets, and an economy showing that these conditions select no limit,
are in \Cref{app:moved-ad}, and the paragraphs this one replaces in
\Cref{app:moved-ah}.

\begin{corollary}[Post-entry non-termination]
\label{cor:nontermination}
Under the hypotheses of \Cref{thm:persistence}, let $T$ be a date with
$x_T^E\in\basin$. For every cutoff $q<\bar x$ and every $n\ge0$,
\[
s_{T+n}^E>q.
\]
Thus every sanction-reduction threshold strictly below the inherited floor
is never reached after basin entry, and when $\bar x>0$ neither is
relief.
\end{corollary}

Relief remains feasible; the corollary says it is not chosen along the
path of the equilibrium $E\in\mathcal E$ after the political state enters the
basin. It is therefore a result about how long sanctions last after basin entry, not about
which disputes enter sanctions in the first place. Only the level condition
is used: at any equilibrium whose sanction rule is at or above the state on
the basin, no cutoff below the inherited floor is reached after entry;
neither A1 nor the target condition enters.

\paragraph{The rule and the closed loop.}
Two abbreviations make the long-run statements readable. Write
\[
R_E(x)=\sigDynTop_{\Vs}\bigl(x,\alpha(x)\bigr),
\qquad
G_E(x)=\F\bigl(x,R_E(x)\bigr),
\]
so that $R_E$ is the sanction rule of $E$ --- the sender's dynamic best
response evaluated at the target's equilibrium policy --- and $G_E$ is the closed
loop it induces on the political state. The realized path of an equilibrium is
then simply $s^E_t=R_E(x^E_t)$ and $x^E_{t+1}=G_E(x^E_t)$; these are the same
objects as in \Cref{sec:model}, named rather than redefined. A \emph{steady
state} of the equilibrium is a fixed point of $G_E$: a political state the closed
loop reproduces. Under strict directional adjustment \textup{(F5)} that is
equivalent to the sanction matching the state there, $R_E(x)=x$, so the two ways
of saying it agree wherever \textup{(F5)} holds.

Three long-run results follow, and they differ in what they assume about the
sanction rule and in how precisely they name the limit. The first
uses only the level condition of \Cref{thm:persistence} --- the sanction rule
at or above the state on the basin --- and establishes convergence to an
unspecified steady state. The second adds monotonicity of the closed loop and
continuity of the rule and pins down the limit as an extremal steady state,
which may differ across starting points. The third retains those conditions but ends the interval at the first
steady state above the pivot, yielding one limit for every start in that interval.

\begin{proposition}[Long-run lock-in level]
\label{prop:longrun}
Let $E\in\mathcal E$, let $\bar x\ge0$, and suppose the sanction rule is at
or above the state on $\basin$, $R_E(x)\ge x$ for $x\in\basin$, with no
floor. With the state evolved by the general law \eqref{eq:transition}
satisfying the maintained \textup{(F1)--(F5)}, the state path from every
basin start converges to some $x^\ast\in\basin$. The realized sanction
sequence converges to the same limit, and $x^\ast$ is an
equilibrium-consistent steady state: sanction $x^\ast$ is optimal at state
$x^\ast$ and reproduces that state. The hypotheses of \Cref{thm:persistence}
imply the level condition.
\end{proposition}

This baseline result requires neither monotonicity nor continuity of the sanction
rule, and neither A1 nor the target condition. Those conditions enter only when the
limit itself is to be pinned down.

Monotonicity of the closed loop and continuity of the sanction rule pin down the limit
(\Cref{prop:order-selection}, \Cref{cor:order-selection-identified} and
\Cref{prop:derived-interval}, stated in \Cref{app:moved-ak}). If the start
lies below the rule's value there, the state converges to the least steady
state of the closed loop at or above it, and if above, to the greatest steady
state at or below it; under strict directional adjustment the realized
sanction converges to the same value. Continuity of the rule is the
requirement, and uniqueness of the maximizer of $J_S$ is one
sufficient condition for it (\Cref{app:moved-ad}). Order selection need not
imply attraction of a whole basin: every start in an invariant interval
converges to one common limit exactly when the interval contains a unique
steady state, and the interval that runs from the basin pivot to the first
steady state above it has that property by construction. In plain terms, the
inherited political state selects, at date zero and within a single
equilibrium, which self-sustaining steady state is reached --- path dependence, not
indeterminacy --- and up to the first such steady state above the pivot every
start reaches it. A steady state of $G_E$ is a political state at which the
sender's optimal sanction reproduces the entrenchment it inherits; none of
the three results delivers a rate (\Cref{app:moved-ah}).

\subsection{Dual-basin path dependence}
\label{sec:mechanism:bistability}

The argument of \Cref{thm:persistence} has a sign-reversed counterpart on a low
political basin, and the two can hold for one equilibrium. With
$\varepsilon=+1$ the high orientation and $\varepsilon=-1$ the low one, the
sender's sanction rule remains non-decreasing in the inherited state, and from
every start in the corresponding basin the state stays there and the
sanction moves monotonically on the basin's side; neither orientation reads
a target condition (\Cref{thm:lowbasin}; the high
orientation reproduces conclusions (i)--(iii) of \Cref{thm:persistence}).
Because the two assumption blocks apply on disjoint state ranges, the same
equilibrium can satisfy both: below a low threshold $\bar x_L$ the state
stays below it with non-increasing sanctions, above a high threshold $\bar x_H$ it
locks in with non-decreasing ones
(\Cref{cor:initial-condition-selection}). This is selection by initial
conditions within one equilibrium, not selection among equilibria; the
appendix supplies an economy satisfying both sets of conditions. The region
between the two thresholds is not left to assumption: if the closed loop
strictly lowers the state on the low basin, strictly raises it on the high
one and has a continuous sanction rule, it determines two of its own steady
states $x_s^-\le x_s^+$ strictly inside $(\bar x_L,\bar x_H)$, with convergence to zero
from every start below $x_s^-$ and lock-in from every start above $x_s^+$,
and the two coincide exactly when there is a single steady state between the
basins (\Cref{prop:derived-dividing-line}; \Cref{fig:bistability}
draws the single-crossing case). The picture --- two
basins inside one equilibrium, with initial conditions selecting between
them --- is the known one of \Cref{sec:literature}; what is new is that the
basins are basins of a political state moved by a coercive instrument.

\subsection{Persistence on a low basin and initial-condition selection}

The persistence argument has a sign-reversed counterpart on a low
political basin. Let $\varepsilon=+1$ denote the high orientation and
$\varepsilon=-1$ the low orientation. In both the sender's sanction rule remains
non-decreasing in the inherited state, and neither reads the target's
sanction preference.

\begin{theorem}[Persistence with orientation $\varepsilon$]
\label{thm:lowbasin}
Fix $\varepsilon\in\{+1,-1\}$ and suppose
A1$_\varepsilon$ holds, together with the orientation-specific level condition;
neither orientation reads a target condition. Suppose
$E=(\Vt,\Vs,\alpha)\in\mathcal E$ with
$\Vt\in\mathcal C_T$, $\Vs\in\mathcal C_S$, and
$\alpha=\alpha_{\Vt,\Vs}$. Then, from every starting state in the
corresponding basin:
\begin{enumerate}[label=(\roman*)]
\item the state remains in that basin;
\item the realized sanction remains on the basin side of its
corresponding extremum;
\item $\varepsilon(s^E_{t+1}-s^E_t)\ge0$ for every $t$.
\end{enumerate}
\end{theorem}

The high orientation reproduces conclusions (i)--(iii) of
\Cref{thm:persistence}; the low orientation keeps the state in the low basin with non-increasing
sanctions. Because the
two assumption blocks apply on disjoint state ranges, the same
equilibrium can satisfy both.

\begin{corollary}[Initial-condition selection under a common equilibrium]
\label{cor:initial-condition-selection}
Suppose the same equilibrium $E$ satisfies the sender's high-basin
conditions with state threshold $\bar x_H$ and sanction threshold
$\bar s_H$ --- at every state of $[\bar x_H,1]$ some maximizer of
$J_S$ is at least $\max\{x,\bar s_H\}$, and the greatest one is
non-decreasing in the state there --- and the sender's low-basin conditions
with thresholds $\bar x_L$ and $\bar s_L$ --- on $[0,\bar x_L]$ the dynamic
best response is at most $\min\{x,\bar s_L\}$ and non-decreasing in the
state --- where $\bar x_L<\bar x_H$. No target condition is used. Then the initial-state regions
$[0,\bar x_L]$ and $[\bar x_H,1]$ are disjoint. For every
$x_0\in[\bar x_H,1]$ and every $t\ge0$,
\[
x_t^E\in[\bar x_H,1],\qquad
s_t^E\ge\max\{\bar x_H,\bar s_H\},\qquad
s_{t+1}^E\ge s_t^E.
\]
For every $x_0\in[0,\bar x_L]$ and every $t\ge0$,
\[
x_t^E\in[0,\bar x_L],\qquad
s_t^E\le\min\{\bar x_L,\bar s_L\},\qquad
s_{t+1}^E\le s_t^E.
\]
\end{corollary}

This is selection by initial conditions within one equilibrium, not
selection among equilibria. The appendix supplies an economy satisfying
both sets of conditions.

\Cref{cor:initial-condition-selection} leaves the region between $\bar x_L$ and
$\bar x_H$ unclassified and takes both basin boundaries as given. Continuity of
the equilibrium closed loop nevertheless derives bounds on the dividing region
without imposing a sign in the middle.

\begin{proposition}[The dividing line is derived]
\label{prop:derived-dividing-line}
Let $E\in\mathcal E$ and $0<\bar x_L\le\bar x_H<1$, for instance the thresholds
of \Cref{cor:initial-condition-selection}. Suppose that the closed loop strictly
lowers the state at every state of $(0,\bar x_L]$ and strictly raises it at
every state of $[\bar x_H,1)$, and that the sanction rule is continuous on
$[0,1]$. The closed loop then maps $[0,1]$ into itself, by feasibility of the
transition and of the rule, and fixes both endpoints of the state space; the
first formalization assumed both. Impose no restriction on the sign
of $G_E(x)-x$ between $\bar x_L$ and $\bar x_H$; monotonicity of the sanction
rule is also unnecessary. Then the equilibrium
closed loop determines two of its own steady states,
\[
  x_s^-=\inf\{\,x\in[\bar x_L,\bar x_H]:x\le G_E(x)\,\},
  \qquad
  x_s^+=\sup\{\,x\in[\bar x_L,\bar x_H]:G_E(x)\le x\,\},
\]
with $\bar x_L<x_s^-\le x_s^+<\bar x_H$, such that every equilibrium path started
anywhere in $[0,x_s^-)$ converges to $0$ and every equilibrium path started
anywhere in $(x_s^+,1]$ converges to $1$. Moreover $x_s^-=x_s^+$ exactly when
$G_E$ has a single steady state in $[\bar x_L,\bar x_H]$.
\end{proposition}

The two convergence regions therefore extend strictly beyond the given basins,
with endpoints derived from the equilibrium; the interval between them need
not collapse to a point, and continuity of the sanction rule is what supplies
a middle steady state (\Cref{app:moved-ah}).

\begin{figure}[tbp]
\centering
\begin{tikzpicture}[
  x=7cm,y=5.2cm,
  >={Stealth[length=2mm]},
  every node/.style={font=\scriptsize}
]
\fill[black!8] (0,0) rectangle (0.22,1);
\fill[black!8] (0.78,0) rectangle (1,1);
\draw[->] (0,0) -- (1.08,0) node[right] {$x_t$};
\draw[->] (0,0) -- (0,1.08) node[above] {$x_{t+1}$};
\draw[black!45,dashed] (0,0) -- (1,1);
\draw[very thick] (0,0) -- (0.22,0.19);
\draw[very thick] (0.78,0.79) -- (1,1);
\draw[black!55,dashed]
  plot[smooth,tension=.75] coordinates
  {(0.22,0.19) (0.38,0.28) (0.50,0.50) (0.66,0.76) (0.78,0.79)};
\draw[->,thick] (0.42,-0.08) -- (0.16,-0.08);
\draw[->,thick] (0.58,-0.08) -- (0.88,-0.08);
\draw[->,thick,dashed] (0.06,-0.19) to[bend right=14]
  node[midway,below] {shock entry when \Cref{thm:shock} applies} (0.97,-0.19);
\node[rotate=90,black!55] at (0.05,0.55) {low basin};
\node[rotate=90,black!55] at (0.975,0.45) {high basin};
\fill (0.5,0.5) circle (0.9pt);
\node[below] at (0.22,0) {$\bar x_L$};
\node[below] at (0.78,0) {$\bar x_H$};
\node[below] at (0.5,0) {$x_s^{\pm}$};
\node[above left,black!55] at (0.62,0.62) {$x_{t+1}=x_t$};
\end{tikzpicture}
\caption{\textbf{Dual-basin selection.}
The solid segments represent the equilibrium dynamics established on the
low and high basins. The dashed middle segment is illustrative. By
\Cref{prop:derived-dividing-line} the middle region contains steady states
$x_s^-\le x_s^+$ derived from the closed loop, with convergence to zero below $x_s^-$ and
lock-in above $x_s^+$; the figure draws the case $x_s^-=x_s^+$ of a single
crossing, which the proposition does not assert. Under
the separate shock-entry conditions, a temporary sanction
shock can move the state into the high basin, after which equilibrium
dynamics resume.}
\label{fig:bistability}
\end{figure}

\subsection{Temporary-shock entry and persistence after entry}
\label{sec:counterfactual}

A temporary sanction shock can move the political state into the high
basin, after which equilibrium policies keep the state there. The
shock is not an equilibrium policy or an optimal sanction policy; each result
describes resumed play under the same specified equilibrium, and every
duration below is a sufficient ceiling that need not be minimal.

\begin{theorem}[Temporary-shock entry and persistence after entry]
\label{thm:shock}
Maintain \textup{(P1)--(P3)} and \textup{(P5)} of \Cref{thm:persistence},
with the same value sets, equilibrium, pivots, and sender Lipschitz constant;
the target condition \textup{(P4)} is not needed. In
addition, suppose the transition follows the general law
\eqref{eq:transition}, satisfies \textup{(F1)--(F4)}, and rises under full
sanctions below the basin, $x<\F(x,1)$ for $x\in[0,\bar x]$ (so $\bar x<1$). Define the full-sanction drift and its minimum below the
basin by
\[
d_{\F}(x):=\F(x,1)-x,
\qquad
\eta_{\F}(\bar x):=\min_{x\in[0,\bar x]}d_{\F}(x).
\]
Then $\eta_{\F}(\bar x)>0$. Define the conservative shock-duration
bound by
\[
\widehat N_{\F}(\bar x)
:=\left\lceil\frac{\bar x}{\eta_{\F}(\bar x)}\right\rceil.
\]
From any $x_0\in[0,1]$, imposing $s_t=1$ for
$t=0,\ldots,\widehat N_{\F}(\bar x)-1$ gives
\[
x_{\widehat N_{\F}(\bar x)}\ge\bar x.
\]
The same conclusion holds for every longer temporary shock at the full sanction.
Consequently $x_{\widehat N_{\F}(\bar x)}\in\basin$, and conclusions
(i)--(iii) of \Cref{thm:persistence} hold along the equilibrium path
resumed from that state.
\end{theorem}

Continuity and full-sanction ascent below the basin make the full-sanction
drift strictly positive on the compact pre-basin interval, where its minimum
is attained. Until the lowest shocked
path enters the basin, it therefore rises by at least
$\eta_{\F}(\bar x)$ each period. The bound is conservative because it
uses only the smallest one-period drift.

The theorem establishes entry followed by persistence. It does not compare
the shocked path with the path the same equilibrium would have followed
unshocked, and so it does not assert that the shock selects a different
long-run outcome; \Cref{prop:activation-not-hysteresis} exhibits an economy
in which it does not. The duration is a sufficient bound that depends on the
economy and is not small in any uniform sense (\Cref{app:moved-ah}).

Under the partial-adjustment law \eqref{eq:affine}, the sharper closed-form duration is
\[
N_*(\bar x,\delta)
:=\left\lceil\frac{\log(1-\bar x)}{\log(1-\delta)}\right\rceil,
\qquad
x_{N_*}=1-(1-\delta)^{N_*}(1-x_0)\ge\bar x.
\]

A transition that carries the state up faster under the full sanction has
a weakly larger minimum entry drift and a weakly shorter sufficient
duration; the drift comparison is strict when the transition comparison is,
while the duration comparison stays weak because distinct drifts can share an
integer ceiling (\Cref{cor:shock-transition-strength}, in \Cref{app:moved-ak}).

\subsection{A temporary shock, then no endogenous exit}
\label{sec:finite-episode}
A confrontation may begin as an attempt to make the other government back
down. The model does not solve that initial interaction, and it carries no
types, beliefs, or signals. What it shows is what a temporary shock
leaves behind. If the shock carries the political state into a region the
resumed equilibrium preserves, the shock's original purpose ceases to
matter for what follows: under unchanged primitives and that equilibrium,
sanctions never fall below the region's floor. This is endogenous non-exit,
not technological impossibility of relief --- the transition still erodes
entrenchment under low sanctions, and \Cref{thm:relief-exit} gives a
sufficient duration for a temporary suspension of sanctions --- and it presupposes that
everything the shock leaves behind is represented by the scalar state.

\begin{theorem}[Temporary shock, then no endogenous exit]
\label{thm:finite-episode}
Maintain \textup{(F1)--(F4)}. Let $(\Vt,\Vs,\alpha)\in\mathcal E$ with
sanction rule $R(x)=\sigDynTop_{\Vs}(x,\alpha(x))$, and let
$\bar x$ and $\bar s$ satisfy $\bar x\le\F(\bar x,\bar s)$ and
$R(x)\ge\bar s$ for every $x\in[\bar x,1]$; both hold whenever the sanction
rule is at or above $\bar x$ on $[\bar x,1]$ and $\bar s=\bar x$. Let
$(u_t)_{t<T}$ be a temporary shock, a feasible sanction path set exogenously before date $T$, with
the equilibrium resumed from date $T$:
\[
x_{t+1}=\F(x_t,s_t),\qquad
s_t=\begin{cases}u_t,&t<T,\\ R(x_t),&t\ge T.\end{cases}
\]
\begin{enumerate}[label=(\roman*)]
\item \emph{Splice.} For $t\ge T$ the path $(x_t,s_t)$ is the equilibrium
path of \Cref{sec:model} restarted at $x_T$.
\item \emph{Non-exit.} If $x_T\in[\bar x,1]$ then $x_t\in[\bar x,1]$ and
$s_t\ge\bar s$ for every $t\ge T$, and for every cutoff $q<\bar s$ there is
no date $t\ge T$ at which $s_t\le q$.
\item \emph{Entry under a non-maximal shock.} Let $0\le\bar x$, let the
sanction level $h$ raise every state of $[0,\bar x]$, $x<\F(x,h)$ there
(held-sanction ascent, which forces $\bar x<h$ and which \textup{(F5)}
implies whenever $\bar x<h\le1$), and let $u_t\ge h$ for every $t<T$. The drift
$\eta_h(\bar x):=\min_{x\in[0,\bar x]}\bigl(\F(x,h)-x\bigr)$ is attained and
strictly positive, and for $T\ge N_h(\bar x):=\lceil\bar x/\eta_h(\bar x)\rceil$
every initial state $x_0\in[0,1]$ gives $x_T\in[\bar x,1]$. At $h=1$,
$N_1(\bar x)$ is $\widehat N_{\F}(\bar x)$ of \Cref{thm:shock}. Where a lower
bound $d>0$ on the transition's sanction slope is available and $\bar x<h\le1$,
$\lceil\bar x/(d(h-\bar x))\rceil$ is an explicit, looser, sufficient
duration.
\item \emph{Ratchet tail.} If $x_T\in[\bar x,1]$ and the sanction rule is at
or above the state on $[\bar x,1]$, the post-shock sanction satisfies
$s_t\ge\max(x_T,\bar x)$ for every $t\ge T$; if the rule is also
non-decreasing on $[\bar x,1]$, the sanction is nondecreasing in $t$. Both
conditions hold under \textup{(P1)--(P5)} of \Cref{thm:persistence} with
pivots $\bar x,\bar s$.
\end{enumerate}
\end{theorem}

The basin conditions of \Cref{thm:finite-episode} are met on the class under
the lifting condition at every equilibrium at which the target maintains full non-compliance, with pivot and
floor both equal to $\underline s=\sigM(0,0)$ ($u_0/q$ on the quadratic--linear family when
$u_0\le q$, as under interiority);
such an equilibrium exists under the entrenchment condition and a non-decreasing
response at non-compliance, which the increasing-differences condition \eqref{eq:rally-margin} implies, the constructed one. That is weak invariance and a floor, not the ratchet and not
hysteresis (\Cref{app:moved-ah}).

The comparison behind (iii), two machine-checked boundaries of the statement
--- a floor equal to the pivot need not be reached in finite time, and no
duration is uniformly small across transitions --- and the instantiation at a
non-affine transition are in \Cref{app:moved-ad}.

\paragraph{The theorem on a finite temporary shock, on the class under the lifting condition.} \emph{(From \Cref{sec:finite-episode}, after \Cref{thm:finite-episode}.)}

The basin conditions of \Cref{thm:finite-episode} are met on the class under
the lifting condition at every equilibrium at which the target maintains full non-compliance, with pivot and
floor both equal to $\underline s=\sigM(0,0)$ ($u_0/q$ on the quadratic--linear family when
$u_0\le q$, as under interiority):
every such equilibrium realizes the constructed response, which is at least
the floor at every state, and $\F(\underline s,\underline s)=\underline s$, so
a temporary shock ending in $[\underline s,1]$ --- in particular one held at
any level above $\underline s$ for the derived duration --- is followed by
permanent residence there with the sanction never below $\underline s$
(\path{placable_post_episode_nonexit_of_policy_zero},
\path{placable_highEpisode_nonexit_of_policy_zero}). Such an equilibrium exists
under the entrenchment condition and a non-decreasing response at non-compliance, which
the increasing-differences condition \eqref{eq:rally-margin} implies, the constructed one; where
none exists the statement is empty. This is weak invariance and a
floor, not the ratchet and not hysteresis: at the affine member the closed
loop contracts and the unshocked path settles above the floor as well.

\subsection{Relief and the exit duration}
\label{sec:relief-exit}

The entry result has a counterpart on the exit side, which uses only the
descent half of \textup{(F5)}, and only inside the basin: the zero sanction
strictly lowers every state of $[\bar x,1]$.
Whether one equilibrium sustains both basins is a further question;
\Cref{prop:two-sided-one-equilibrium} is the one economy at which it is settled
(\Cref{app:moved-ah}).

\begin{theorem}[Exit under a temporary suspension of sanctions]
\label{thm:relief-exit}
Maintain \textup{(F1)--(F4)}, let $\bar x\le1$, and suppose the zero
sanction strictly lowers every state of the basin, $\F(x,0)<x$ for
$x\in[\bar x,1]$ (post-basin descent, which forces $\bar x>0$ and which
\textup{(F5)} implies whenever $\bar x>0$). Define the relief
drift and its minimum inside the basin by
\[
  r_{\F}(x):=x-\F(x,0),
  \qquad
  \eta^{R}_{\F}(\bar x):=\min_{x\in[\bar x,1]}r_{\F}(x).
\]
The minimum is attained and $\eta^{R}_{\F}(\bar x)>0$. Let
\[
  \widehat N^{R}_{\F}(\bar x)
  :=\left\lceil\frac{1-\bar x}{\eta^{R}_{\F}(\bar x)}\right\rceil .
\]
If the sanction is held at zero from an arbitrary feasible initial state
$x_0\in[0,1]$, then the resulting state path satisfies $x_n\le\bar x$ at every
date $n\ge\widehat N^{R}_{\F}(\bar x)$.
\end{theorem}

As with \Cref{thm:shock}, the relief path is exogenous, not chosen: it is a
temporary shock during which neither government optimizes, and the result is
correspondingly not an optimal-relief claim. What it adds is the second
threshold. Under the hypotheses of \Cref{thm:persistence} relief never arrives
endogenously --- inside the basin the equilibrium sanction stays above the state
inherited at entry, so the set of dates at which the equilibrium grants relief
is empty --- and under the weaker floor of \Cref{thm:finite-episode} it never
falls below the region's floor. Read together, the two statements say something
neither says alone:

\begin{quote}
The equilibrium never ends the sanctions regime on its own. A temporary
suspension of sanctions of the stated sufficient duration crosses the threshold;
whether sanctions then stay low depends on a separate low-basin
condition for the resumed equilibrium.
\end{quote}

The bound is a sufficient duration for a temporary suspension
of sanctions, not a minimal one. What this implies for sunset clauses and review
dates, and how the entry and exit durations compare, is in
\Cref{app:moved-ah,app:moved-ad}.

Two limits should be stated plainly. \Cref{thm:relief-exit} stops at the
crossing: it establishes that the state re-enters $[0,\bar x]$ and by itself
says nothing about whether it \emph{remains} there once the shock ends. That
gap is closed \emph{conditionally}, with the same shape as the entry side:

Conditional on the low-basin level condition --- the sanction rule at most
the state and at most a level $\underline s$ on $[0,\bar x]$ --- a
temporary suspension of sanctions of the stated length is followed by a state path that
stays in $[0,\bar x]$ and is non-increasing (\Cref{prop:post-relief-persistence}). The level condition is the dual of the one
\Cref{thm:shock} assumes on the high basin: an assumption about the
equilibrium in question, discharged at the dual-basin economy of
\Cref{prop:two-sided-one-equilibrium} and at no other economy in the
development. Read with \Cref{thm:shock}, the two statements are two
opposite-direction crossings, each followed under its own conditions by the state staying on its new side; they
are not a same-equilibrium hysteresis unless both level conditions hold for
one equilibrium. Two further limits, at the threshold itself and in the
conservatism of the duration, are in \Cref{app:moved-ah}.

Relief in this section is set exogenously. At the constructed forward-looking
equilibrium of the class under the lifting condition it does not arise endogenously, since the
target maintains non-compliance there; against a myopic target, and below the intermediate range of the target's discount factor in \Cref{prop:bracket},
compliance buys it (\Cref{tab:four-cells}). Under the entrenchment condition and
a non-decreasing response, at every equilibrium of the class at which the
target maintains full non-compliance (the constructed one among them) a spell of
zero sanctions of any length, followed by that equilibrium, is worth
strictly less to the sender than the equilibrium itself --- by at least a
stage-gain floor per period ($u_0^2/(2q)$ on the quadratic--linear family), when
no compliance level is worth more to the sender than non-compliance at zero sanction ---
and weakly less to the target, whose entrenchment
condition makes the discounted entrenchment a sanction builds cover its pain;
their gains sum to a negative number, so no lump-sum side payment between them
makes the spell acceptable to both (\Cref{prop:relief-participation}). The
obstacle to de-escalation there is participation before it is enforcement.
Outside the entrenchment condition the target's half fails: a temporary suspension
of sanctions can benefit a target that would still choose non-compliance.

At the dual-basin economy there is one equilibrium at which a temporary
full-sanction shock of the stated length is followed by permanent residence
in the high basin with the sanction at its maximum, and a temporary suspension
of sanctions by permanent residence in the low basin with a zero
sanction (\Cref{prop:two-sided-one-equilibrium} in \Cref{app:dual-basin-economy}).
This is the same-equilibrium two-sided statement that \Cref{sec:relief-exit}
declines to make in general. It is affine, it leaves the interval between the
two basins unclassified, and it is an existence certificate
(\Cref{app:moved-ah}).


\paragraph{Forward-looking incentives that change behavior do not imply hysteresis.}
\label{sec:compstat:activation-not-hysteresis}
On the class under the lifting condition (\Cref{sec:existence:placable}) both forward-looking incentives
change behavior at a constructed equilibrium, and one might read that as history
dependence. It is not, and the most transparent member of its quadratic--linear
family shows why.

\begin{proposition}[A contracting closed loop]
\label{prop:activation-not-hysteresis}
On every economy in the class under the lifting condition let $\rho$ be the
sender's dynamic best response at full non-compliance, evaluated at its constructed value, which
exists on every such economy; under the entrenchment condition and the increasing-differences condition
\eqref{eq:rally-margin} the map $G(x)=\F(x,\rho(x))$ is the state law of the
equilibrium of
\Cref{prop:placable-family}. Then $G$ is Lipschitz on $[0,1]$ with constant
$L_{\mathrm{loop}}=g_{\mathrm{hi}}+d_{\mathrm{hi}}\cdot L_{\mathrm{resp}}/\mu$,
where $\mu=q-\beta_S(\sigma_V d_{\mathrm{hi}}^2+K_S M_s)$ is the
strong-concavity modulus of the sender's dynamic objective, $M_s$ the
sanction curvature bound of the transition, and
$L_{\mathrm{resp}}=r_{\mathrm{hi}}+\beta_S\bigl(\bar\sigma g_{\mathrm{hi}}d_{\mathrm{hi}}
+(K_S+2\bar\sigma)M_{\mathrm{mix}}\bigr)$ with
$\bar\sigma=\max(\sigma_V,\sigma_C)$ and $K_S=w_{\mathrm{hi}}/(1-\beta_Sg_{\mathrm{hi}})$. Two paths that resume the equilibrium
at states $x$ and $y$ are within $L_{\mathrm{loop}}^n|x-y|$ after $n$
periods. On every member of the relief class with $L_{\mathrm{loop}}<1$ that
satisfies myopic interiority $u_{\mathrm{hi}}<q$, the entrenchment condition and
the increasing-differences condition, the equilibrium has exactly one steady state in $[0,1]$, at
which the sanction equals the state, and every initial state converges to
it. At the affine member $K_S=2$,
$\mu=23/8$, $L_{\mathrm{resp}}=5/8$, and $L_{\mathrm{loop}}=14/23$; at the
non-affine member $K_S=15/7$, $\mu=141/50$, $L_{\mathrm{resp}}=17/25$, and
$L_{\mathrm{loop}}=35/47$.
\end{proposition}

So at both members both governments are forward-looking, both forward-looking incentives change
behavior, and a temporary sanction shock has no asymptotic effect, though
transient and cumulative effects are not excluded; the absence of an
asymptotic effect holds on the whole relief class wherever
$L_{\mathrm{loop}}<1$ and the stated conditions hold. The persistence in
\Cref{thm:shock} rests on its basin conditions, not on whether the forward-looking incentives
change behavior. The two constants are evaluated upper bounds, not measured
convergence rates (\Cref{app:moved-ah}).

% Moved from the main text, 2026-10-16 (structure plan S2): the reading of lock-in and the
% three notions of stability, formerly in Section 5.2.
\subsection{Persistence after relief}

\begin{proposition}[Post-relief persistence, conditional on the low-basin
level condition]
\label{prop:post-relief-persistence}
Maintain \textup{(F1)--(F4)}, let $\bar x\le1$, suppose the zero sanction
strictly lowers every state of $[\bar x,1]$, $\F(x,0)<x$ there (which forces
$\bar x>0$), and let $R$ be the sanction rule of an equilibrium's
value and policy. Suppose the low-basin
level condition holds: for some $\underline s$, $R(x)\le\min(x,\underline s)$
on $[0,\bar x]$. Then from every initial state, after a temporary suspension
of sanctions of length $\widehat N^{R}_{\F}(\bar x)$, the state path evolved by
$\F$ under $R$ lies in $[0,\bar x]$ at every later date and is
non-increasing.
\end{proposition}

The level condition is the dual of the one \Cref{thm:shock} assumes on the
high basin: it is an assumption about the equilibrium in question, not a
property of the model, and it is discharged at the dual-basin economy of
\Cref{prop:two-sided-one-equilibrium} (where the rule is zero on the low
basin) and at no other economy in the development. Read with
\Cref{thm:shock}, the two statements are two opposite-direction crossings
with conditional persistence on each side; they are not a same-equilibrium
hysteresis unless both level conditions hold for one equilibrium, which is
what \Cref{prop:two-sided-one-equilibrium} exhibits. Two further limits, at the
threshold itself and in the conservatism of the duration, are in
\Cref{app:moved-ah}.

\subsection{Reading lock-in}
\label{app:persistence:reading}

Whether persistent sanctions are permanent is a question about the
dynamics, which lock-in answers equilibrium by equilibrium. Lock-in is a property of a given equilibrium: its
direction and level conditions are imposed only at basin states, at every value
pair in the stated sets, and the level condition only against the state itself
(\Cref{prop:lock-in}). These requirements can be met: on a class of corner
economies, defined by inequalities among primitives, at which the constructed
equilibrium sanctions fully at every state, and at an example economy whose
continuation values change observable choices.
Lock-in's comparative statics run in one direction: with the same transition and
initial state, a non-decreasing sanction rule at least as high as another on the
high basin yields a state path at least as high at every date.

Three notions of stability appear in these results, and they remain
distinct. First, state paths converge. When the equilibrium law of motion is
non-decreasing on a basin that it preserves and the sanction rule is continuous
there, the start fixes the steady state reached: the least one at or above the
start if the sanction there is at least the state, and otherwise the greatest one
at or below it. Second, equilibrium uniqueness holds relative to stated
restrictions: invariant sets of value functions, tie-breaking rules and
non-decreasing sanction rules (\Cref{sec:results:security}). Uniqueness in
dynamic games of this kind is open \citep{DoraszelskiSatterthwaite2010}. On every
family the paper constructs, the sender's optimal sanction is unique at each state
and feasible compliance level, so the greatest, the least and any other selection
coincide there. Third, in a stochastic extension of the model studied elsewhere in this
appendix, an invariant distribution of the state exists; its analytic
existence argument and its formalized constructions have different verification
scopes.
